Department of Computer Science\\
New Jersey Institute of Technology\\
\texttt{ea442@njit.edu}\\
Supervisor: Dr. Martin Kellogg\\[0.8em] % MK: I don't use my middle initial in professional settings, as some people do. Note for example that all of my papers list me as "Martin Kellogg", while Mike's list him as "Michael D. Ernst". In general, you should assume that this kind of choice is intentional and try to match whatever the person themselves usually does (it's a politeness thing and some people will take offense, although I am not offended).

\textbf{Abstract}

\todo{Erfan: reconcile the following two paragraph, which resulted from a failed git merge. The reason that I usually put many line breaks into paper text is specifically
to avoid this problem, and since you haven't done so with this abstract, I am going to let you deal with the fallout. Perhaps that will encourage you to add more line breaks
in the future. (Recall that git is line-based, so if edits don't touch the same lines they won't conflict. In paper text with line breaks, I always put todo comments on new
lines, guaranteeing the absence of conflict.)}

\todo{This version of the paragraph is Martin's.} Developers spend a significant portion of their time understanding source code, yet code comprehensibility remains\todo{``is?'' Since you haven't brought up other studies yet, it doesn't make sense to use ``remains''} difficult to define and measure. Existing studies rely heavily on indirect proxies (e.g., time, accuracy, or syntactic metrics), many of which have not been systematically validated. In prior work\todo{You should not refer to the ASE 2026 submission as ``prior work'', since it forms a part of your thesis. Instead, use the present tense: ``We address this gap...'' It's important to make sure that the committee understands the ASE 2026 study as your work, and referring to it as ``prior work'' might give the mistaken impression that you weren't involved. It is appropriate to refer to e.g., Oscar and I's FSE 2023 paper as ``prior work'', but not something that this document will describe in complete detail.}, we addressed this gap by introducing a practical conceptualization of code comprehensibility based on expert agreement and constructing ground-truth rankings using a Delphi study. We then evaluated commonly used proxies against this ground truth, showing that behavior-based proxies—particularly those based on input--output reasoning—align more closely with expert judgments, while syntactic proxies often exhibit weak or even negative correlations.

\todo{This version of the paragraph is Erfan's.} Developers spend a significant portion of their time understanding source code, yet code comprehensibility remains difficult to define and measure. Existing studies rely heavily on indirect proxies (e.g., time spent producing an output for a given input), many of which have not been systematically validated. In prior work, we addressed this gap by introducing a practical conceptualization of code comprehensibility based on expert agreement and constructing ground-truth rankings using a Delphi study. We then evaluated commonly used proxies against this ground truth, showing that proxies derived from input–output questions and those based on response time rather than accuracy are particularly reliable, whereas proxies derived from questions about program syntax (rather than semantics) are especially unreliable, regardless of the measurement strategy.


Building on these findings, this proposal investigates whether reliable automatic proxies for code comprehensibility can be developed and used to support scalable analysis and tooling\todo{too vague---scalable analysis and tooling \emph{for what?} I think you want to just be a bit more direct about your goal here: automatically estimating comprehensibility without having to ask developers.}. In particular, we propose a controlled experimental study to examine the relationship between code verifiability and comprehensibility, testing the hypothesis that code that is easier to verify is also easier to understand. Beyond this study, we outline a broader research agenda that includes validating proxies under varying levels of accidental complexity, reassessing prior work in light of proxy reliability, exploring LLM-based agents as substitutes for human participants, and developing tools that leverage reliable proxies to improve code quality. Together, this work aims to establish a principled and empirically grounded foundation for measuring and improving code comprehensibility.
